% Options for packages loaded elsewhere
\PassOptionsToPackage{unicode}{hyperref}
\PassOptionsToPackage{hyphens}{url}
%
\documentclass[
  10pt,
]{article}
\usepackage{amsmath,amssymb}
\usepackage{iftex}
\ifPDFTeX
  \usepackage[T1]{fontenc}
  \usepackage[utf8]{inputenc}
  \usepackage{textcomp} % provide euro and other symbols
\else % if luatex or xetex
  \usepackage{unicode-math} % this also loads fontspec
  \defaultfontfeatures{Scale=MatchLowercase}
  \defaultfontfeatures[\rmfamily]{Ligatures=TeX,Scale=1}
\fi
\usepackage{lmodern}
\ifPDFTeX\else
  % xetex/luatex font selection
\fi
% Use upquote if available, for straight quotes in verbatim environments
\IfFileExists{upquote.sty}{\usepackage{upquote}}{}
\IfFileExists{microtype.sty}{% use microtype if available
  \usepackage[]{microtype}
  \UseMicrotypeSet[protrusion]{basicmath} % disable protrusion for tt fonts
}{}
\makeatletter
\@ifundefined{KOMAClassName}{% if non-KOMA class
  \IfFileExists{parskip.sty}{%
    \usepackage{parskip}
  }{% else
    \setlength{\parindent}{0pt}
    \setlength{\parskip}{6pt plus 2pt minus 1pt}}
}{% if KOMA class
  \KOMAoptions{parskip=half}}
\makeatother
\usepackage{xcolor}
\usepackage[margin=0.8in]{geometry}
\usepackage{color}
\usepackage{fancyvrb}
\newcommand{\VerbBar}{|}
\newcommand{\VERB}{\Verb[commandchars=\\\{\}]}
\DefineVerbatimEnvironment{Highlighting}{Verbatim}{commandchars=\\\{\}}
% Add ',fontsize=\small' for more characters per line
\newenvironment{Shaded}{}{}
\newcommand{\AlertTok}[1]{\textcolor[rgb]{1.00,0.00,0.00}{\textbf{#1}}}
\newcommand{\AnnotationTok}[1]{\textcolor[rgb]{0.38,0.63,0.69}{\textbf{\textit{#1}}}}
\newcommand{\AttributeTok}[1]{\textcolor[rgb]{0.49,0.56,0.16}{#1}}
\newcommand{\BaseNTok}[1]{\textcolor[rgb]{0.25,0.63,0.44}{#1}}
\newcommand{\BuiltInTok}[1]{\textcolor[rgb]{0.00,0.50,0.00}{#1}}
\newcommand{\CharTok}[1]{\textcolor[rgb]{0.25,0.44,0.63}{#1}}
\newcommand{\CommentTok}[1]{\textcolor[rgb]{0.38,0.63,0.69}{\textit{#1}}}
\newcommand{\CommentVarTok}[1]{\textcolor[rgb]{0.38,0.63,0.69}{\textbf{\textit{#1}}}}
\newcommand{\ConstantTok}[1]{\textcolor[rgb]{0.53,0.00,0.00}{#1}}
\newcommand{\ControlFlowTok}[1]{\textcolor[rgb]{0.00,0.44,0.13}{\textbf{#1}}}
\newcommand{\DataTypeTok}[1]{\textcolor[rgb]{0.56,0.13,0.00}{#1}}
\newcommand{\DecValTok}[1]{\textcolor[rgb]{0.25,0.63,0.44}{#1}}
\newcommand{\DocumentationTok}[1]{\textcolor[rgb]{0.73,0.13,0.13}{\textit{#1}}}
\newcommand{\ErrorTok}[1]{\textcolor[rgb]{1.00,0.00,0.00}{\textbf{#1}}}
\newcommand{\ExtensionTok}[1]{#1}
\newcommand{\FloatTok}[1]{\textcolor[rgb]{0.25,0.63,0.44}{#1}}
\newcommand{\FunctionTok}[1]{\textcolor[rgb]{0.02,0.16,0.49}{#1}}
\newcommand{\ImportTok}[1]{\textcolor[rgb]{0.00,0.50,0.00}{\textbf{#1}}}
\newcommand{\InformationTok}[1]{\textcolor[rgb]{0.38,0.63,0.69}{\textbf{\textit{#1}}}}
\newcommand{\KeywordTok}[1]{\textcolor[rgb]{0.00,0.44,0.13}{\textbf{#1}}}
\newcommand{\NormalTok}[1]{#1}
\newcommand{\OperatorTok}[1]{\textcolor[rgb]{0.40,0.40,0.40}{#1}}
\newcommand{\OtherTok}[1]{\textcolor[rgb]{0.00,0.44,0.13}{#1}}
\newcommand{\PreprocessorTok}[1]{\textcolor[rgb]{0.74,0.48,0.00}{#1}}
\newcommand{\RegionMarkerTok}[1]{#1}
\newcommand{\SpecialCharTok}[1]{\textcolor[rgb]{0.25,0.44,0.63}{#1}}
\newcommand{\SpecialStringTok}[1]{\textcolor[rgb]{0.73,0.40,0.53}{#1}}
\newcommand{\StringTok}[1]{\textcolor[rgb]{0.25,0.44,0.63}{#1}}
\newcommand{\VariableTok}[1]{\textcolor[rgb]{0.10,0.09,0.49}{#1}}
\newcommand{\VerbatimStringTok}[1]{\textcolor[rgb]{0.25,0.44,0.63}{#1}}
\newcommand{\WarningTok}[1]{\textcolor[rgb]{0.38,0.63,0.69}{\textbf{\textit{#1}}}}
\usepackage{longtable,booktabs,array}
\usepackage{calc} % for calculating minipage widths
% Correct order of tables after \paragraph or \subparagraph
\usepackage{etoolbox}
\makeatletter
\patchcmd\longtable{\par}{\if@noskipsec\mbox{}\fi\par}{}{}
\makeatother
% Allow footnotes in longtable head/foot
\IfFileExists{footnotehyper.sty}{\usepackage{footnotehyper}}{\usepackage{footnote}}
\makesavenoteenv{longtable}
\setlength{\emergencystretch}{3em} % prevent overfull lines
\providecommand{\tightlist}{%
  \setlength{\itemsep}{0pt}\setlength{\parskip}{0pt}}
\setcounter{secnumdepth}{5}
\ifLuaTeX
\usepackage[bidi=basic]{babel}
\else
\usepackage[bidi=default]{babel}
\fi
\babelprovide[main,import]{english}
% get rid of language-specific shorthands (see #6817):
\let\LanguageShortHands\languageshorthands
\def\languageshorthands#1{}
\usepackage{xurl}
\usepackage{microtype}
\setlength{\emergencystretch}{3em}
\ifLuaTeX
  \usepackage{selnolig}  % disable illegal ligatures
\fi
\usepackage[]{natbib}
\bibliographystyle{plainnat}
\usepackage{bookmark}
\IfFileExists{xurl.sty}{\usepackage{xurl}}{} % add URL line breaks if available
\urlstyle{same}
\hypersetup{
  pdftitle={Broadband Antennas Integrated with Radome Faces for Passive Detection of Clandestine or Anomalous RF Signals},
  pdfauthor={Literature review based on Consensus search results},
  pdflang={en},
  hidelinks,
  pdfcreator={LaTeX via pandoc}}

\title{Broadband Antennas Integrated with Radome Faces for Passive
Detection of Clandestine or Anomalous RF Signals}
\author{Literature review based on Consensus search results}
\date{6 August 2026}

\begin{document}
\maketitle

{
\setcounter{tocdepth}{3}
\tableofcontents
}
\section{Search method}\label{search-method}

This review used a \textbf{Decomposition framework} with five sub-areas:

\begin{enumerate}
\def\labelenumi{\arabic{enumi}.}
\tightlist
\item
  broadband antenna architectures and radome-face placement;
\item
  electromagnetic effects and functional radomes;
\item
  passive emitter detection, direction finding, and localization;
\item
  clandestine, anomalous, and deviated-signal identification; and
\item
  receiver integration, calibration, and operational validation.
\end{enumerate}

The search comprised one reconnaissance query followed by a ten-query
standard review: five sub-area searches, two review searches, two
era-gated searches (through 2015 and from 2021 onward), and one
follow-up centered on the most-cited paper. The available Consensus
connector returned at most ten papers per query. Filters such as
\texttt{human} and minimum sample size were inapplicable to this
engineering topic; the connector did not expose SJR filtering. Only
papers returned and fetched from Consensus during the session are cited
here.

\section{Topic overview}\label{topic-overview}

The literature supports the feasibility of integrating broadband or
ultra-wideband receiving antennas with conformal radomes for passive RF
monitoring, but the evidence is divided into largely separate streams.
The \textbf{antenna/radome literature} is comparatively mature, with
measured prototypes evaluating impedance bandwidth, gain, insertion
loss, angular stability, cavity resonances, and radar cross-section. The
\textbf{passive detection literature} is also mature at the algorithmic
level, particularly for spectrum sensing, direction of arrival (DOA),
time difference of arrival (TDOA), frequency difference of arrival
(FDOA), and specific-emitter identification. The sparse area is their
end-to-end intersection: few returned studies jointly demonstrate a
broadband conformal radome aperture, a calibrated multichannel receiver,
detection of previously unknown emissions, emitter classification, and
localization under realistic platform and propagation conditions. In
particular, stealth-oriented frequency-selective-surface (FSS) radomes
may deliberately reject out-of-band signals, creating a direct conflict
with the objective of listening over the broadest possible spectrum
\citep{costa2012, lv2021, sheng2024, xing2022}.

The closest terminology for the combined topic is: \textbf{conformal
ultra-wideband radio direction finder}, \textbf{passive emitter
surveillance}, \textbf{ESM/ELINT receiving array},
\textbf{radome-integrated spectrum-monitoring aperture},
\textbf{specific emitter identification}, \textbf{open-set RF anomaly
detection}, and \textbf{co-aperture passive direction-finding antenna}.

A crucial systems distinction is that \textbf{antenna coverage
bandwidth} is not equivalent to \textbf{instantaneous receiver
bandwidth}. A physically broadband antenna may span many octaves while
the RF front end, ADC, channelizer, data transport, or processor
observes only a much narrower portion at any instant
\citep{flak2022, madanayake2024, subbaraman2022}.

\section{Start here: priority reading
order}\label{start-here-priority-reading-order}

\subsection{1. Closest review to the complete antenna
problem}\label{closest-review-to-the-complete-antenna-problem}

\href{https://consensus.app/papers/details/04301748949257caa17485d3480352f2/?utm_source=chatgpt}{\textbf{Wideband
Antennas of Passive Seekers for Anti Radiation Missiles}}, Gupta et
al.~(2023), compares spiral, log-periodic, printed Vivaldi, and
all-metal Vivaldi antennas and distinguishes concealed-behind-radome,
flush-mounted, and conformal arrangements \citep{gupta2023}. It is the
most direct returned overview of antenna placement for passive wideband
reception. While reading, check whether its application-specific
assumptions on size, frequency range, polarization, and angular coverage
match the intended platform.

\subsection{2. Seminal functional-radome
paper}\label{seminal-functional-radome-paper}

\href{https://consensus.app/papers/details/1811b78b36fa5cc0840c16f9dd6f2a1b/?utm_source=chatgpt}{\textbf{A
Frequency Selective Radome With Wideband Absorbing Properties}}, Costa
and Monorchio (2012), establishes an influential architecture combining
an FSS passband with wideband out-of-band absorption \citep{costa2012}.
The key design tension is that a radome advantageous for stealth may be
inappropriate for broad-spectrum interception because out-of-band
absorption removes potentially relevant emissions.

\subsection{3. Foundational integrated conformal direction
finder}\label{foundational-integrated-conformal-direction-finder}

\href{https://consensus.app/papers/details/5f166a5e2c715f25a2cfa7d55f20749f/?utm_source=chatgpt}{\textbf{Design
and Full-Wave Analysis of Conformal Ultra-Wideband Radio Direction
Finders}}, Caratelli, Liberal, and Yarovoy (2011), develops a conformal
circular array covering 250 MHz--3.3 GHz and explicitly includes
coupling compensation, array geometry, and ambiguity analysis
\citep{caratelli2011}. Its main lesson is that calibration and manifold
distortion are as important as nominal antenna bandwidth for reliable
DOA estimation.

\subsection{4. Broad orientation for emitter
identity}\label{broad-orientation-for-emitter-identity}

\href{https://consensus.app/papers/details/85a1acc05180584992b4bd752a62599e/?utm_source=chatgpt}{\textbf{A
Comprehensive Survey on Radio Frequency Fingerprinting: Traditional
Approaches, Deep Learning, and Open Challenges}}, Jagannath, Jagannath,
and Kumar (2022), connects SIGINT, RF fingerprinting, conventional
features, deep learning, datasets, and security applications
\citep{jagannath2022}. Pay particular attention to the difference
between closed-set device classification and recognizing an emitter that
was absent from training.

\subsection{5. Current radome frontier}\label{current-radome-frontier}

\href{https://consensus.app/papers/details/6cb5b6043f755325b7abd988656c9afd/?utm_source=chatgpt}{\textbf{A
Conformal Miniaturized Bandpass Frequency-Selective Surface With Stable
Frequency Response for Radome Applications}}, Sheng et al.~(2024),
reports measured stability for a curved X-band FSS, including large
incidence angles and bending \citep{sheng2024}. Frequency stability
alone, however, does not prove preservation of phase, group delay, or
array-manifold accuracy required for direction finding.

\subsection{6. Current open-set detection
frontier}\label{current-open-set-detection-frontier}

\href{https://consensus.app/papers/details/66ce93d9ddd05433ac370321a0d58745/?utm_source=chatgpt}{\textbf{An
Open-Set Supervised Anomaly Detection Method for Unauthorized
Broadcasting Identification}}, Zhang et al.~(2025), directly addresses
unauthorized broadcasting in an open-set setting using temporal
representation learning and support-vector data description
\citep{zhang2025open}. The central evaluation issue is whether
performance remains stable across receivers, propagation channels,
center-frequency changes, and truly unseen emitters.

\subsection{7. Key controversy and
vulnerability}\label{key-controversy-and-vulnerability}

\href{https://consensus.app/papers/details/29ea0ff4b51857cca284e60d8c9c9d19/?utm_source=chatgpt}{\textbf{Robustness
of Deep Learning-Based Specific Emitter Identification under Adversarial
Attacks}}, Sun et al.~(2022), shows that small deliberately constructed
perturbations can seriously degrade deep-learning-based specific-emitter
identification, while adversarial training partially restores
performance \citep{sun2022}. A classifier with high ordinary test
accuracy is therefore not automatically suitable for detecting a
deliberately camouflaged or deviated signal.

\section{Cross-search signals}\label{cross-search-signals}

\subsection{Repeat-hit papers}\label{repeat-hit-papers}

No paper met the strict criterion of appearing in three distinct
sub-area searches. This is itself informative: the five branches remain
weakly integrated.

The strongest repeat hits across the reconnaissance, sub-area, and
review stages were:

\begin{itemize}
\tightlist
\item
  Gupta et al.~(2023): three appearances overall \citep{gupta2023}.
\item
  Gharat et al.~(2026): three appearances overall \citep{gharat2026}.
\item
  Costa and Monorchio (2012): functional-radome search and seminal-paper
  follow-up \citep{costa2012}.
\item
  Soltanieh et al.~(2020): anomaly/SEI search and review search
  \citep{soltanieh2020}.
\item
  Xing et al.~(2022): radome search and post-2021 era search
  \citep{xing2022}.
\item
  Madanayake et al.~(2024): integration search and post-2021 era search
  \citep{madanayake2024}.
\item
  Tahseen et al.~(2021): reconnaissance and radome-review search
  \citep{tahseen2021}.
\end{itemize}

These repetitions identify useful bridge papers, but they should not be
interpreted as proof of foundational status by themselves.

\subsection{Citation velocity}\label{citation-velocity}

Citation velocity was approximated as the Consensus citation count
divided by \texttt{max(1,\ 2026\ -\ publication\ year)}. It is
descriptive only: recent survey articles are favored, and database
citation counts can change.

\begin{longtable}[]{@{}lrr@{}}
\toprule\noalign{}
Paper & Consensus citations & Approx. citations/year \\
\midrule\noalign{}
\endhead
\bottomrule\noalign{}
\endlastfoot
Jagannath et al.~2022 \citep{jagannath2022} & 263 & 65.8 \\
Soltanieh et al.~2020 \citep{soltanieh2020} & 339 & 56.5 \\
Costa and Monorchio 2012 \citep{costa2012} & 678 & 48.4 \\
Taherpour et al.~2010 \citep{taherpour2010} & 441 & 27.6 \\
Sheng et al.~2024 \citep{sheng2024} & 53 & 26.5 \\
Lv et al.~2021 \citep{lv2021} & 67 & 13.4 \\
Xing et al.~2022 \citep{xing2022} & 47 & 11.8 \\
Omar and Shen 2018 \citep{omar2018} & 77 & 9.6 \\
\end{longtable}

Normalized citation rate favors recent surveys, while total citations
more clearly identify Costa--Monorchio and Taherpour et al.~as
long-standing reference points.

\section{How the field got here}\label{how-the-field-got-here}

The early literature treated passive RF direction finding largely as an
array-design and estimation problem. By 2010--2012, researchers were
integrating broadband conformal arrays with calibration methods and
developing blind multiple-antenna spectrum detectors
\citep{caratelli2011, liberal2011, taherpour2010}. During the same
period, frequency-selective radomes evolved from protective dielectric
covers into functional electromagnetic structures capable of
transmitting selected bands while absorbing others \citep{costa2012}.

The next phase emphasized flush integration, cavity control, compact FSS
structures, and low-RCS antenna--radome co-design
\citep{lv2021, omar2018, tianang2017, xing2022}. In parallel,
\textbf{specific emitter identification} increasingly became known as
\textbf{RF fingerprinting}, moving from handcrafted transient and
spectrum features toward deep learning and open-set anomaly detection
\citep{jagannath2022, soltanieh2020, zhang2025open}. Recent work adds
real-time channelization, FPGA/SDR acceleration, and adaptive
monitoring, but the hardware and signal-intelligence branches remain
only partially integrated
\citep{flak2022, gharat2026, madanayake2024, sorecau2026}.

\begin{longtable}[]{@{}
  >{\raggedleft\arraybackslash}p{(\columnwidth - 4\tabcolsep) * \real{0.4000}}
  >{\raggedright\arraybackslash}p{(\columnwidth - 4\tabcolsep) * \real{0.3000}}
  >{\raggedright\arraybackslash}p{(\columnwidth - 4\tabcolsep) * \real{0.3000}}@{}}
\toprule\noalign{}
\begin{minipage}[b]{\linewidth}\raggedleft
Year
\end{minipage} & \begin{minipage}[b]{\linewidth}\raggedright
Milestone
\end{minipage} & \begin{minipage}[b]{\linewidth}\raggedright
Significance
\end{minipage} \\
\midrule\noalign{}
\endhead
\bottomrule\noalign{}
\endlastfoot
2010 & Blind multiple-antenna spectrum sensing \citep{taherpour2010} &
Establishes detection under unknown channel, noise, and source
parameters. \\
2011 & Conformal UWB radio direction finders
\citep{caratelli2011, liberal2011} & Integrates wideband elements,
circular geometry, coupling analysis, and calibration. \\
2012 & Wideband absorptive frequency-selective radome \citep{costa2012}
& Makes the radome an active part of spectral selectivity and signature
control. \\
2012 & Compact 70--3000 MHz DF antenna \citep{bailey2012} & Demonstrates
end-to-end receiver/antenna integration over a broad band. \\
2017--2018 & Flush Vivaldi and thin conformal 3-D FSS systems
\citep{tianang2017, omar2018} & Addresses cavity resonances, structural
thickness, and curved integration. \\
2020--2022 & RF fingerprinting surveys and deep-learning transition
\citep{soltanieh2020, jagannath2022} & Shifts identification from
handcrafted features toward data-driven models. \\
2021--2024 & Hybrid radomes and co-designed antenna--FSS structures
\citep{lv2021, sheng2024, xing2022} & Combines passband transmission,
absorption, diffusion, conformality, and RCS control. \\
2024--2026 & Multichannel spectrum intelligence and open-set anomaly
monitoring
\citep{gharat2026, madanayake2024, sorecau2026, zhang2025open} & Moves
toward real-time detection, directional channelization, and
unknown-emitter recognition. \\
\end{longtable}

\subsection{Terminology shifts}\label{terminology-shifts}

\begin{itemize}
\tightlist
\item
  \textbf{Radio direction finder / RDF} became \textbf{DOA estimation},
  \textbf{passive emitter localization}, and \textbf{spectrum
  intelligence}.
\item
  \textbf{Emitter recognition} became \textbf{specific emitter
  identification / SEI}, then \textbf{RF fingerprinting / RFF}.
\item
  \textbf{Illegal radio monitoring} increasingly appears as
  \textbf{unauthorized broadcasting identification} and \textbf{open-set
  RF anomaly detection}.
\item
  \textbf{Frequency-selective radome} expanded into \textbf{absorptive
  FSR}, \textbf{rasorber radome}, \textbf{absorptive-diffusive radome},
  and \textbf{active/tunable FSS}.
\item
  \textbf{Shared aperture} also appears as \textbf{co-aperture},
  \textbf{common-radome}, and \textbf{multifunctional aperture}.
\end{itemize}

\section{Sub-area guide 1: broadband antenna architectures and
radome-face
placement}\label{sub-area-guide-1-broadband-antenna-architectures-and-radome-face-placement}

\subsection{What the research shows}\label{what-the-research-shows}

Conformal elliptical dipoles, spirals, and Vivaldi/tapered-slot antennas
are recurring candidates because they combine broad impedance bandwidth
with geometries suitable for circular, cylindrical, or flush
installation
\citep{caratelli2011, gupta2023, liberal2011, pi2025, tianang2017}.
Wideband DOA performance depends strongly on the calibrated array
manifold: mutual coupling, structural scattering, element curvature, and
radome effects can introduce angular bias even when return loss and gain
appear acceptable \citep{caratelli2011, liberal2011}.

Cavity-backed Vivaldi structures can deliver flush mounting, but cavity
eigenmodes may produce severe in-band resonances unless explicitly
suppressed \citep{tianang2017}. A recent co-aperture spiral concept
proposes 2--40 GHz coverage under a common radome, but the returned work
reports simulation rather than complete field validation \citep{pi2025}.

\subsection{Key papers}\label{key-papers}

\begin{itemize}
\tightlist
\item
  \href{https://consensus.app/papers/details/04301748949257caa17485d3480352f2/?utm_source=chatgpt}{\textbf{Wideband
  Antennas of Passive Seekers for Anti Radiation Missiles}}, Gupta et
  al.~(2023), Consensus citation count 1: best returned comparison of
  concealed, flush, and conformal passive-seeker arrangements
  \citep{gupta2023}.
\item
  \href{https://consensus.app/papers/details/5f166a5e2c715f25a2cfa7d55f20749f/?utm_source=chatgpt}{\textbf{Design
  and Full-Wave Analysis of Conformal Ultra-Wideband Radio Direction
  Finders}}, Caratelli et al.~(2011), 17 citations: integrates element,
  array, calibration, and DOA considerations \citep{caratelli2011}.
\item
  \href{https://consensus.app/papers/details/962901757e715c669ba9962b91c0b96b/?utm_source=chatgpt}{\textbf{Conformal
  Antenna Array for Ultra-Wideband Direction-of-Arrival Estimation}},
  Liberal et al.~(2011), 9 citations: explicitly includes a protective
  radome and coupling compensation \citep{liberal2011}.
\item
  \href{https://consensus.app/papers/details/747aac0a449e5797ad83ea444fb20680/?utm_source=chatgpt}{\textbf{Ultra-Wideband
  Lossless Cavity-Backed Vivaldi Antenna}}, Tianang et al.~(2017), 42
  citations: important treatment of cavity-induced resonances in flush
  arrays \citep{tianang2017}.
\item
  \href{https://consensus.app/papers/details/7848c93584d354f19ce7554e9d95de54/?utm_source=chatgpt}{\textbf{A
  Compact Airborne Co-Aperture Cavity-Backed Spiral Antenna Design}}, Pi
  (2025), 0 citations: emerging common-radome concept spanning two bands
  over 2--40 GHz \citep{pi2025}.
\end{itemize}

\subsection{Search terms}\label{search-terms}

conformal UWB antenna; passive seeker antenna; radio direction finder;
cavity-backed Vivaldi; cavity-backed spiral; sinuous antenna;
flush-mounted broadband array; co-aperture antenna; amplitude-tracking
antenna set; phase-matched DF antenna.

\subsection{Boolean searches}\label{boolean-searches}

\begin{Shaded}
\begin{Highlighting}[]
\NormalTok{("conformal antenna" OR "flush{-}mounted antenna" OR "embedded antenna")}
\NormalTok{AND ("ultra{-}wideband" OR wideband)}
\NormalTok{AND ("direction finding" OR "passive emitter" OR ESM OR ELINT)}
\end{Highlighting}
\end{Shaded}

\begin{Shaded}
\begin{Highlighting}[]
\NormalTok{(Vivaldi OR "tapered slot" OR spiral OR sinuous OR "log{-}periodic")}
\NormalTok{AND (radome OR "common aperture" OR "co{-}aperture")}
\NormalTok{AND (receive OR surveillance OR "direction finding")}
\end{Highlighting}
\end{Shaded}

\section{Sub-area guide 2: electromagnetic effects and functional radome
technologies}\label{sub-area-guide-2-electromagnetic-effects-and-functional-radome-technologies}

\subsection{What the research shows}\label{what-the-research-shows-1}

The radome must be treated as part of the receiving transfer function
rather than merely as a mechanical cover. FSS and hybrid radomes can
provide low-loss transmission windows together with out-of-band
absorption or diffuse scattering \citep{costa2012, lv2021, xing2022}.
Curved and conformal implementations increasingly demonstrate stable
center frequency under oblique incidence, but most report insertion loss
and radiation performance rather than phase-delay uniformity across an
array aperture \citep{omar2018, sheng2024}. A recent conformal-FSS
mini-review also identifies materials, simulation, manufacturing, and
practical implementation as continuing challenges \citep{korkut2024}.

For broad-spectrum monitoring, the design objective differs from that of
a conventional stealth radome. A narrow FSS passband may protect an
intended radar or communications band while making the system deaf to
unauthorized emissions outside that window
\citep{costa2012, tahseen2021}. The radome therefore needs either a
genuinely broad transparent region, multiple transparent windows, or a
controllable transmission state.

\subsection{Key papers}\label{key-papers-1}

\begin{itemize}
\tightlist
\item
  \href{https://consensus.app/papers/details/1811b78b36fa5cc0840c16f9dd6f2a1b/?utm_source=chatgpt}{\textbf{A
  Frequency Selective Radome With Wideband Absorbing Properties}}, Costa
  and Monorchio (2012), 678 citations: seminal passband-plus-absorber
  architecture \citep{costa2012}.
\item
  \href{https://consensus.app/papers/details/905ccf75368255a78a236f083a31e9ee/?utm_source=chatgpt}{\textbf{Thin
  3-D Bandpass Frequency-Selective Structure Based on Folded Substrate
  for Conformal Radome Applications}}, Omar and Shen (2018), 77
  citations: fabricated semicylindrical FSS radome with broadband-horn
  integration \citep{omar2018}.
\item
  \href{https://consensus.app/papers/details/62fa6b42c0305286b6a3a4c60ea60a46/?utm_source=chatgpt}{\textbf{Hybrid
  Absorptive-Diffusive Frequency Selective Radome}}, Lv et al.~(2021),
  67 citations: broad transmission window with absorption and diffusion
  sidebands \citep{lv2021}.
\item
  \href{https://consensus.app/papers/details/7a9db905b6b953c3a6f429fe051d4096/?utm_source=chatgpt}{\textbf{A
  Low-RCS and Wideband Circularly Polarized Array Antenna Co-Designed
  With a High-Performance AMC-FSS Radome}}, Xing et al.~(2022), 47
  citations: measured co-design with less than 1 dB reported gain
  deterioration across its operating band \citep{xing2022}.
\item
  \href{https://consensus.app/papers/details/6cb5b6043f755325b7abd988656c9afd/?utm_source=chatgpt}{\textbf{A
  Conformal Miniaturized Bandpass Frequency-Selective Surface With
  Stable Frequency Response for Radome Applications}}, Sheng et
  al.~(2024), 53 citations: strong recent result on curved-FSS angular
  stability \citep{sheng2024}.
\end{itemize}

\subsection{Search terms}\label{search-terms-1}

frequency-selective radome; FSS radome; rasorber; absorptive
frequency-selective radome; transmitarray radome; radome insertion loss;
radome boresight error; radome depolarization; radome group delay;
active tunable radome.

\subsection{Boolean searches}\label{boolean-searches-1}

\begin{Shaded}
\begin{Highlighting}[]
\NormalTok{(radome OR "antenna cover")}
\NormalTok{AND ("frequency selective surface" OR FSS OR rasorber OR metasurface)}
\NormalTok{AND (wideband OR broadband OR multiband)}
\NormalTok{AND ("insertion loss" OR "angular stability" OR "boresight error")}
\end{Highlighting}
\end{Shaded}

\begin{Shaded}
\begin{Highlighting}[]
\NormalTok{("conformal radome" OR "curved FSS")}
\NormalTok{AND (transmission OR passband)}
\NormalTok{AND (phase OR "group delay" OR depolarization OR aberration)}
\NormalTok{AND (array OR "direction finding")}
\end{Highlighting}
\end{Shaded}

\section{Sub-area guide 3: passive emitter detection, direction finding,
and
localization}\label{sub-area-guide-3-passive-emitter-detection-direction-finding-and-localization}

\subsection{What the research shows}\label{what-the-research-shows-2}

Blind multiple-antenna spectrum sensing can outperform conventional
energy detection when channel gain, source statistics, and noise
variance are uncertain \citep{taherpour2010}. Once a signal is detected,
localization commonly uses AoA/DOA, TDOA, FDOA, or combinations of them.
TDOA/FDOA methods can achieve strong accuracy, but receiver geometry,
synchronization, platform motion, and computational load are central
constraints \citep{liu2019, pine2021}.

For radome-mounted arrays, calibration is not an optional
post-processing refinement. Directional errors can create biased
bearings or ghost emitter positions; joint localization and
self-calibration methods attempt to correct those effects
\citep{zhang2023calibration}. The older conformal-RDF literature already
recognized mutual coupling and support-structure distortion as
array-manifold problems \citep{caratelli2011, liberal2011}.

\subsection{Key papers}\label{key-papers-2}

\begin{itemize}
\tightlist
\item
  \href{https://consensus.app/papers/details/dd361298468d58a9a3f6b6af0c34e977/?utm_source=chatgpt}{\textbf{Multiple
  Antenna Spectrum Sensing in Cognitive Radios}}, Taherpour et
  al.~(2010), 441 citations: foundational blind multiple-antenna
  detector under parameter uncertainty \citep{taherpour2010}.
\item
  \href{https://consensus.app/papers/details/5f166a5e2c715f25a2cfa7d55f20749f/?utm_source=chatgpt}{\textbf{Design
  and Full-Wave Analysis of Conformal Ultra-Wideband Radio Direction
  Finders}}, Caratelli et al.~(2011), 17 citations: connects antenna
  nonidealities directly to DOA performance \citep{caratelli2011}.
\item
  \href{https://consensus.app/papers/details/47ae35140d8e563486d4808c6acd8485/?utm_source=chatgpt}{\textbf{Computationally
  Efficient TDOA and FDOA Estimation Algorithm in Passive Emitter
  Localisation}}, Liu et al.~(2019), 14 citations: targets real-time
  feasibility and multi-emitter cross-term suppression \citep{liu2019}.
\item
  \href{https://consensus.app/papers/details/035e2d1ded6e5316b49b04c2007445cb/?utm_source=chatgpt}{\textbf{The
  Geometry of Far-Field Passive Source Localization With TDOA and
  FDOA}}, Pine et al.~(2021), 36 citations: clarifies when TDOA/FDOA
  measurements geometrically constrain far-field sources
  \citep{pine2021}.
\item
  \href{https://consensus.app/papers/details/a785d6a4988c51bcb3d4c3b896010b09/?utm_source=chatgpt}{\textbf{Passive
  Joint Emitter Localization with Sensor Self-Calibration}}, Zhang et
  al.~(2023), 8 citations: addresses calibration bias in distributed
  passive arrays \citep{zhang2023calibration}.
\end{itemize}

\subsection{Search terms}\label{search-terms-2}

passive emitter localization; radio direction finding; angle of arrival;
direction of arrival; amplitude-comparison DF; phase interferometer;
TDOA/FDOA; direct position determination; array-manifold calibration;
noncooperative emitter.

\subsection{Boolean searches}\label{boolean-searches-2}

\begin{Shaded}
\begin{Highlighting}[]
\NormalTok{("passive emitter" OR "noncooperative emitter" OR "radio source")}
\NormalTok{AND (localization OR geolocation OR "direction finding")}
\NormalTok{AND (DOA OR AOA OR TDOA OR FDOA OR interferometry)}
\end{Highlighting}
\end{Shaded}

\begin{Shaded}
\begin{Highlighting}[]
\NormalTok{("wideband direction finding" OR "broadband direction finding")}
\NormalTok{AND (calibration OR "array manifold" OR coupling OR radome)}
\NormalTok{AND (multichannel OR array)}
\end{Highlighting}
\end{Shaded}

\section{Sub-area guide 4: clandestine, anomalous, and deviated-signal
identification}\label{sub-area-guide-4-clandestine-anomalous-and-deviated-signal-identification}

\subsection{What the research shows}\label{what-the-research-shows-3}

RF fingerprinting uses hardware-dependent imperfections to associate a
received waveform with a device or emitter. Surveys show a progression
from manually designed transient and modulation features to deep neural
representations \citep{jagannath2022, soltanieh2020}. Much of the
literature, however, assumes a \textbf{closed set}, in which every
emitter observed at test time belongs to a known training class.

Clandestine-signal detection usually requires an open-set or anomaly
formulation. Illegal-FM monitoring has been demonstrated using
specific-emitter features and outlier detection \citep{guo2021}, while
recent work explicitly models unauthorized transmissions not fully
represented in training \citep{zhang2025open}. Frequency-hopping
emitters can also be fingerprinted \citep{kang2021}, but adversarial
perturbations may make deep SEI systems unreliable against an adaptive
opponent \citep{sun2022}.

\subsection{Key papers}\label{key-papers-3}

\begin{itemize}
\tightlist
\item
  \href{https://consensus.app/papers/details/1c3cd42586205e0280356e90cb7ab147/?utm_source=chatgpt}{\textbf{A
  Review of Radio Frequency Fingerprinting Techniques}}, Soltanieh et
  al.~(2020), 339 citations: core taxonomy of transient and
  modulated-signal fingerprints \citep{soltanieh2020}.
\item
  \href{https://consensus.app/papers/details/85a1acc05180584992b4bd752a62599e/?utm_source=chatgpt}{\textbf{A
  Comprehensive Survey on Radio Frequency Fingerprinting: Traditional
  Approaches, Deep Learning, and Open Challenges}}, Jagannath et
  al.~(2022), 263 citations: broad connection among SIGINT, datasets,
  ML, and RFF \citep{jagannath2022}.
\item
  \href{https://consensus.app/papers/details/b81586a8d9ee5753a4b48df60568f585/?utm_source=chatgpt}{\textbf{An
  SEI-Based Identification Scheme for Illegal FM Broadcast}}, Guo et
  al.~(2021), 7 citations: direct application to legal-versus-illegal FM
  emitters \citep{guo2021}.
\item
  \href{https://consensus.app/papers/details/29ea0ff4b51857cca284e60d8c9c9d19/?utm_source=chatgpt}{\textbf{Robustness
  of Deep Learning-Based Specific Emitter Identification under
  Adversarial Attacks}}, Sun et al.~(2022), 22 citations: demonstrates
  an important security weakness in neural SEI \citep{sun2022}.
\item
  \href{https://consensus.app/papers/details/66ce93d9ddd05433ac370321a0d58745/?utm_source=chatgpt}{\textbf{An
  Open-Set Supervised Anomaly Detection Method for Unauthorized
  Broadcasting Identification}}, Zhang et al.~(2025), 0 citations:
  closest direct match to unknown unauthorized-signal detection
  \citep{zhang2025open}.
\end{itemize}

\subsection{Search terms}\label{search-terms-3}

specific emitter identification; SEI; radio-frequency fingerprinting; RF
fingerprint; unauthorized broadcasting; illegal-emitter identification;
open-set RF classification; unknown emitter; spectrum anomaly detection;
low-probability-of-intercept signal; frequency-hopping emitter
identification.

\subsection{Boolean searches}\label{boolean-searches-3}

\begin{Shaded}
\begin{Highlighting}[]
\NormalTok{("specific emitter identification" OR "RF fingerprinting" OR "radio frequency fingerprint")}
\NormalTok{AND ("open set" OR unknown OR unauthorized OR illegal OR anomaly)}
\end{Highlighting}
\end{Shaded}

\begin{Shaded}
\begin{Highlighting}[]
\NormalTok{("clandestine signal" OR "unauthorized transmission" OR "illegal broadcast"}
\NormalTok{ OR "deceptive emitter" OR spoofing)}
\NormalTok{AND (detection OR identification OR classification)}
\NormalTok{AND (spectrum OR RF OR radio)}
\end{Highlighting}
\end{Shaded}

\begin{Shaded}
\begin{Highlighting}[]
\NormalTok{("frequency hopping" OR "low probability of intercept" OR burst)}
\NormalTok{AND ("emitter identification" OR "anomaly detection" OR fingerprinting)}
\end{Highlighting}
\end{Shaded}

\section{Sub-area guide 5: integration, calibration, and operational
validation}\label{sub-area-guide-5-integration-calibration-and-operational-validation}

\subsection{What the research shows}\label{what-the-research-shows-4}

Antenna bandwidth does not determine how much spectrum can be observed
simultaneously. SDR systems remain limited by ADC rate, RF-front-end
linearity, host-transfer bandwidth, channelization complexity, storage,
and processing \citep{flak2022, subbaraman2022}. Fast sweeping increases
frequency coverage but introduces blind time, creating a risk of missing
short bursts or hopping signals.

Recent systems address this through FPGA acceleration, multichannel
SDRs, directional sub-band processing, and adaptive high-resolution
analysis \citep{gharat2026, madanayake2024, sorecau2026}. Receiver
saturation and intermodulation are especially important for
clandestine-signal detection because a weak unauthorized emitter may
coexist with much stronger legitimate transmitters. The fetched
Consensus record for the recent multipath RFSoC front end did not expose
its abstract, so only its bibliographic metadata is used here
\citep{kovacs2025}.

\subsection{Key papers}\label{key-papers-4}

\begin{itemize}
\tightlist
\item
  \href{https://consensus.app/papers/details/fe357d36f1dc5626b1f9d7fae0ff75a7/?utm_source=chatgpt}{\textbf{Hardware-Accelerated
  Real-Time Spectrum Analyzer With a Broadband Fast Sweep Feature Based
  on the Cost-Effective SDR Platform}}, Flak (2022), 16 citations:
  quantifies real-time bandwidth and sweep-speed trade-offs
  \citep{flak2022}.
\item
  \href{https://consensus.app/papers/details/f527168041d95e7ba00d1be0e2d79a21/?utm_source=chatgpt}{\textbf{Observing
  Wideband RF Spectrum with Low-Cost, Resource-Limited SDRs}},
  Subbaraman et al.~(2022), 0 citations: identifies backhaul and
  processing as barriers to real-time wideband observation
  \citep{subbaraman2022}.
\item
  \href{https://consensus.app/papers/details/1f33f743e5f75b38b47a1905f27c7a5f/?utm_source=chatgpt}{\textbf{Design
  of Multichannel Spectrum Intelligence Systems Using Approximate DFT
  for Antenna-Array Spectrum Perception}}, Madanayake et al.~(2024), 13
  citations: integrates directional sensing, channelization, and
  adaptive spectral attention \citep{madanayake2024}.
\item
  \href{https://consensus.app/papers/details/9d218ae20f9b5c1dae9b17a3583aece9/?utm_source=chatgpt}{\textbf{Wideband
  Monitoring System of Drone Emissions Based on SDR Technology with
  RFNoC Architecture}}, Sorecau et al.~(2026), 3 citations: multichannel
  architecture intended to capture frequency-hopping drone emissions
  \citep{sorecau2026}.
\item
  \href{https://consensus.app/papers/details/6719301371565fb4bacbf2c123a0d12d/?utm_source=chatgpt}{\textbf{Real-Time
  Passive RF Spectrum Monitoring and Localization Using Adaptive Signal
  Analysis}}, Gharat et al.~(2026), 0 citations: combines adaptive
  anomaly thresholds and directional localization in a low-cost
  framework \citep{gharat2026}.
\end{itemize}

\subsection{Search terms}\label{search-terms-4}

instantaneous RF bandwidth; wideband spectrum monitoring; SDR
channelizer; RFSoC spectrum analyzer; multichannel digital receiver;
receiver blind time; dynamic range; spurious-free dynamic range; array
calibration; real-time spectrum intelligence.

\subsection{Boolean searches}\label{boolean-searches-4}

\begin{Shaded}
\begin{Highlighting}[]
\NormalTok{("wideband spectrum monitoring" OR "spectrum intelligence")}
\NormalTok{AND (SDR OR RFSoC OR FPGA OR channelizer)}
\NormalTok{AND ("real time" OR multichannel)}
\end{Highlighting}
\end{Shaded}

\begin{Shaded}
\begin{Highlighting}[]
\NormalTok{("passive RF monitoring" OR "emitter surveillance")}
\NormalTok{AND ("dynamic range" OR linearity OR saturation OR "blind time")}
\NormalTok{AND (array OR antenna OR receiver)}
\end{Highlighting}
\end{Shaded}

\section{Key research groups}\label{key-research-groups}

The returned Consensus metadata did not expose affiliations. Therefore,
affiliations are deliberately not inferred or invented; they should be
checked on the linked paper pages or publisher records.

\subsection{I. Liberal, D. Caratelli, and A.
Yarovoy}\label{i.-liberal-d.-caratelli-and-a.-yarovoy}

Their recurring work covers conformal UWB arrays, radio direction
finding, coupling compensation, and array-manifold calibration. Three
closely related pre-2015 records appeared in the era-gated search. A
representative paper is Caratelli et al.~(2011) \citep{caratelli2011}.

\subsection{F. Costa, A. Monorchio, and G.
Manara}\label{f.-costa-a.-monorchio-and-g.-manara}

This group recurs around frequency-selective surfaces, electromagnetic
absorbers, equivalent-circuit modeling, and functional radomes. The
representative paper in the returned set is Costa and Monorchio (2012)
\citep{costa2012}.

\subsection{Zhongxiang Shen and
collaborators}\label{zhongxiang-shen-and-collaborators}

The returned papers cover conformal 3-D FSS structures,
absorptive-diffusive radomes, wide-angle transmission, and RCS
reduction. Recurring collaborators include A. Omar, Qihao Lv, and C.
Jin. A representative paper is Lv et al.~(2021) \citep{lv2021}.

\subsection{X. Sheng, Ning Liu, and
collaborators}\label{x.-sheng-ning-liu-and-collaborators}

Their work centers on equivalent-circuit-based bandpass FSS design,
conformal stability, wide rejection bands, and radome applications. A
representative paper is Sheng et al.~(2024) \citep{sheng2024}.

\subsection{M. Sorecau, E. Sorecau, P. Bechet, and
collaborators}\label{m.-sorecau-e.-sorecau-p.-bechet-and-collaborators}

The returned work focuses on SDR spectrum monitoring, swept versus
real-time observation, multichannel RFNoC architectures, and
drone-emission detection. A representative paper is Sorecau et
al.~(2026) \citep{sorecau2026}.

\section{Open questions and gaps}\label{open-questions-and-gaps}

\subsection{Methodological gaps}\label{methodological-gaps}

\subsubsection{Lack of end-to-end
evaluation}\label{lack-of-end-to-end-evaluation}

Most antenna and radome papers terminate at S-parameters, gain,
radiation patterns, RCS, or transmission coefficients. Detection papers
generally begin with already digitized signals. Few returned studies
measured the full chain

\begin{quote}
radome -\textgreater{} antenna array -\textgreater{} RF front end
-\textgreater{} ADC/channelizer -\textgreater{} detector -\textgreater{}
classifier -\textgreater{} direction/localization output
\end{quote}

under the same experimental conditions.

\textbf{Why it matters:} a radome may have acceptable average insertion
loss while introducing frequency- and angle-dependent phase errors that
degrade DOA or RF fingerprints.

\subsubsection{Closed-set and dataset-leakage
risks}\label{closed-set-and-dataset-leakage-risks}

RF-fingerprinting studies often evaluate known devices under conditions
similar to training \citep{jagannath2022, soltanieh2020}. Open-set work
is emerging \citep{zhang2025open}, but channel, receiver, and location
effects can accidentally become part of the learned fingerprint.

\textbf{Why it matters:} a system may classify the laboratory channel or
receiver chain rather than the clandestine transmitter.

\subsubsection{Insufficient adversarial
testing}\label{insufficient-adversarial-testing}

Sun et al.~show that deliberately perturbed emissions can defeat deep
SEI \citep{sun2022}. This issue is not routinely incorporated into
antenna/radome or spectrum-monitoring validation.

\textbf{Why it matters:} a clandestine source may intentionally imitate
an authorized waveform or suppress its hardware fingerprint.

\subsubsection{Inconsistent metrics}\label{inconsistent-metrics}

The branches use different outcome measures: VSWR, gain, insertion loss,
and RCS for hardware; probability of detection and false alarm for
sensing; angular RMSE for localization; and accuracy or AUC for
classification.

\textbf{Why it matters:} there is no accepted system-level metric
connecting radome loss and distortion to the probability of detecting
and identifying a weak emitter.

\subsection{Population and context
gaps}\label{population-and-context-gaps}

\subsubsection{No demonstrated seamless multi-octave
system}\label{no-demonstrated-seamless-multi-octave-system}

The returned works cover useful but different ranges: 250 MHz--3.3 GHz
\citep{caratelli2011, liberal2011}, 1.5--7.5 GHz \citep{tianang2017},
X-band radomes \citep{sheng2024, tahseen2021}, and a proposed 2--40 GHz
co-aperture receiver \citep{pi2025}.

\textbf{Why it matters:} a single broad-spectrum radome face may require
several sub-arrays and receiver paths rather than one universal element.

\subsubsection{Weak, short-duration, and frequency-hopping
emitters}\label{weak-short-duration-and-frequency-hopping-emitters}

Real-time monitoring papers recognize hopping or transient signals
\citep{flak2022, sorecau2026}, but complete evaluation with weak
emitters, strong adjacent transmitters, multipath, and platform motion
is sparse.

\textbf{Why it matters:} sweep-based receivers can miss short
transmissions, while wide instantaneous receivers may saturate or exceed
processing capacity.

\subsubsection{Environmental and structural
aging}\label{environmental-and-structural-aging}

The returned conformal-radome studies focus mostly on electromagnetic
prototypes. Thermal gradients, moisture, vibration, erosion,
manufacturing tolerances, and long-term dielectric changes are not
integrated with classification and DOA calibration.

\textbf{Why it matters:} slow structural drift can look like a signal
deviation or corrupt a stored emitter baseline.

\subsubsection{Airborne and curved-platform field
datasets}\label{airborne-and-curved-platform-field-datasets}

Several papers target airborne or UAV platforms
\citep{gupta2023, pi2025, sorecau2026}, but public, reproducible
datasets combining curved arrays, radome distortion, and emitter
identity appear sparse in the returned literature.

\textbf{Why it matters:} laboratory models may not transfer to a moving
platform with changing orientation, polarization, and multipath.

\subsection{Conceptual and theoretical
gaps}\label{conceptual-and-theoretical-gaps}

\subsubsection{Stealth radome versus listening
radome}\label{stealth-radome-versus-listening-radome}

FSS radomes commonly suppress out-of-band exposure
\citep{costa2012, lv2021, xing2022}. A surveillance system, by contrast,
seeks sensitivity to signals outside the expected band.

\textbf{Why it matters:} optimizing for low RCS and optimizing for
broad-spectrum listening can be contradictory objectives. A tunable,
segmented, or deliberately wide-transmission radome may be required.

\subsubsection{Emitter identity versus signal
deviation}\label{emitter-identity-versus-signal-deviation}

RF fingerprinting asks, ``Which physical transmitter produced this
waveform?'' Anomaly detection asks, ``Does this observation differ from
expected behavior?'' These are related but not equivalent.

\textbf{Why it matters:} a known transmitter can change modulation or
operating parameters without changing hardware, while a deceptive
transmitter can imitate the nominal waveform but retain a different
hardware fingerprint.

\subsubsection{Antenna bandwidth versus receiver
observability}\label{antenna-bandwidth-versus-receiver-observability}

A 2--40 GHz antenna does not imply simultaneous 38 GHz observation.
Receiver bandwidth, channelization, and data transport determine what is
actually visible \citep{flak2022, madanayake2024, subbaraman2022}.

\textbf{Why it matters:} system specifications should separately state
antenna coverage, instantaneous bandwidth, sweep revisit time, and
probability of intercept.

\subsubsection{Propagation fingerprint versus emitter
fingerprint}\label{propagation-fingerprint-versus-emitter-fingerprint}

Radome response, antenna pattern, multipath, receiver nonlinearity, and
transmitter hardware all affect the measured waveform.

\textbf{Why it matters:} without domain-invariant models or controlled
calibration, a classifier can confuse location and channel changes with
a new or deviating emitter.

\section{Recommended research
architecture}\label{recommended-research-architecture}

The literature supports a modular research program rather than
attempting a universal antenna immediately:

\begin{enumerate}
\def\labelenumi{\arabic{enumi}.}
\tightlist
\item
  Segment the desired spectrum into technically realistic sub-bands.
\item
  Use conformal or flush broadband elements appropriate to each band,
  such as spiral or sinuous structures at lower bands and Vivaldi or
  tapered-slot elements at higher bands.
\item
  Treat the radome as a calibrated complex transfer function
  \(H(f,\theta,\phi,p)\), including amplitude, phase, polarization, and
  incidence angle.
\item
  Use synchronized multichannel reception for DOA and a separate
  wideband or swept path for discovery.
\item
  Separate the inference chain into signal-presence detection, waveform
  characterization, known-emitter matching, open-set anomaly detection,
  and localization.
\item
  Evaluate the complete system using both ordinary and deliberately
  deceptive emitters while varying SNR, channel, orientation,
  temperature, and radome condition.
\end{enumerate}

The most defensible thesis-level gap emerging from the returned
literature is:

\begin{quote}
\textbf{Joint co-design and experimental validation of a conformal
broadband antenna--radome--receiver system for open-set passive emitter
detection, specific-emitter identification, and direction finding, with
explicit calibration of radome- and platform-induced distortions.}
\end{quote}

\section{Limitations of this review}\label{limitations-of-this-review}

The review is bounded by the papers returned by Consensus in this
session, the ten-results-per-query limit, and the metadata exposed by
the connector. Citation counts are provisional and will change. Some
technically relevant defense and industrial work may be classified,
proprietary, poorly indexed, or published outside peer-reviewed venues.
The exact-title follow-up for the most-cited radome paper returned only
four results; no outside material was silently added.

\renewcommand\refname{References}
  \bibliography{references}

\end{document}
